\chapter{Discussion}
\label{sec:discussion}

\section{Evaluation}

As mentioned in Section~\ref{sec:proof-tactics-map},
the automatic proof tactic is limited at this stage.
Therefore, it is premature to create an extensive benchmark suite for testing.
Nonetheless, we can evaluate the implementation using several manually created example tests,
as well as modified existing tests from the Boogie and Proofgen repositories that involve maps.
Most working examples are interger or boolean maps,
that use map select and store operations,
but do not involve axioms with quantifications.

Here is an example we are currently not able to automatcally certify:

\begin{lstlisting}[caption={Non-working test}]
function dummy<T>(x: T) : T; //to force Boogie not to monomorphize

axiom (forall i:int, m:[int]int :: i > 0 ==> m[i] == 1);
axiom (forall i:int, m:[int]int :: i <= 0 ==> m[i] == -1);

procedure p ()
{
    var m: [int] int;
    var x: int;
    assert m[x]*m[x] == 1;
}
\end{lstlisting}

\subsection{Newly Supported Boogie Tests}

The prior work \cite{parthasarathy2021formally} included seven external benchmarks \cite{chen2017triggerless}, that initially used maps.
They rewrote it to avoid using maps, and used it as additional tests.
We used the original tests, rewrote them to avoid the convenient map assign operator,
and run our tool.
Five of these tests failed. Mostly due to proof tactics failing.
Some of the tests also require type helper lemmas, that are not yet implemented in the semantics.
Two managed to verify, \texttt{sum\_of\_arrary.bpl} and \texttt{max\_of\_arrary\_v1.bpl}.
Both are similar, so we will show one of them:

\begin{lstlisting}[caption={\texttt{sum\_of\_array.bpl}}]
function dummy<T>(x: T) : T; //to force Boogie not to monomorphize

// sum of content of a from position to position
function sum_from_to([int]int, int, int) returns(int);
axiom ( forall a: [int]int, low, high: int ::
    high < low   ==>  sum_from_to(a,low,high) == 0 );
axiom ( forall a: [int]int, low, high: int :: high >= low
    ==> sum_from_to(a,low,high-1) + a[high]
        == sum_from_to(a,low,high) );

procedure sum (a: [int]int, n: int) returns(s: int)
	requires n >= 0;
	ensures s == sum_from_to(a, 1, n);
{
	var i: int;
	i := 1;  s := 0;
	while ( i <= n )
	invariant ( 1 <= i  &&  i <=  n + 1);
	invariant ( s == sum_from_to(a, 1, i - 1) );
	{
		s := s + a[i];
		i := i + 1;
	}
}
\end{lstlisting}

This shows a more elaborate example,
with axioms and loop invariants, although only reading from maps.


\subsection{Dafny Tests}
In Section~\ref{sec:motivation}, one motivation is to eventually be able to support Boogie files, generated by Dafny.
In Section~\ref{sec:unsupported-usage} we was, that even the simplest Dafny program
generates a large Boogie program.
Unfortunately, even after having some support for maps is not enough to verify it.
There are more features missing for that.
Currently it fails, as \lstinline[columns=fixed]{free} procedures are not supported by Proofgen.
So one need to either continue working on the next unsupported feature,
or implement a pruning operation, as not all generated Boogie code is necessary for the Dafny input.


\subsection{Prior Tests}
We also ran the existing test suite of the proof generation, that use no maps.
We want to see, if all our changes broke some existing tests.
Fortunately, out of the 120 tests that were passing before our changes,
only one (\texttt{Lock.bpl}) is failing after our change.
Investigating this the failure revealed, that there is an issues in the final step,
with quantification and type parameters.
This seems like a similar issue described in Section~\ref{sec:adjust-foundational},
but we did not manage to find the proper fix for it yet.


\subsection{VC change Tests}
In Section~\ref{sec:our-vc-change} we described, that we modified the VC,
by weakening one axiom.
We want to see if this change alone leads to less verification by Boogie
(without the proof generating module).
So we tested it with the Boogie and Dafny test suit and compared it
with a Boogie version without this change:

[TODO rephrase] The prior work \cite{parthasarathy2021formally} included seven external benchmarks \cite{chen2017triggerless}, that initially used maps.

[TODO evaluation still to be done]


\section{Limitations}
\label{sec:limitations}

\subsection{Program-Dependent Semantics}
Currently, the implementation is fixed at three levels of map nesting.
While this serves as a practical limitation for the prototype,
it is theoretically extensible to any finite depth required by a target Boogie program.
This depth must be determined in advance so that we can appropriately define the types and functions of the map semantics.
Consequently, different input programs may yield distinct instantiations of the semantics.
A notable characteristic of this approach is that the value semantics inherently become program-dependent,
implying that a single, fully uniform semantic definition may not cover all programs.


\subsection{VC Change}
One necessary modification was to adjust the VC map axiom,
as it was too general for our semantics.
For the cases we tested, this change made no difference.
In general, the modified axiom is not weaker for type-safe programs.
Nevertheless, it may affect the SMT solver's ability to check the VC.


\subsection{Single-Arity Maps}
Boogie allows maps to take multiple keys (see Section~\ref{sec:semantics-boogie-maps}).
Our semantics, and therefore the proof generation, are currently limited to a single key.
As previously mentioned, we can systematically transform multi-arity maps into single-arity maps via currying,
where the range is itself a map that takes the subsequent key.
We adopted this limitation to make the initial project scope more feasible.
Early conceptual designs relied on this assumption,
as it helped ensure that the domain used only lower levels
and that either the domain or range resided at its maximal potential level.

Using the well-formed type definition and inductive structure described in Section~\ref{sec:wf-definition},
it should be possible to natively support multi-arity maps.
This would require changing the domain representation to a list of types,
while \isa{FunL} would take a list of values as its domain.
The crucial adjustment in \isa{wf{\isacharunderscore}{\kern0pt}L} would be ensuring that \emph{all} domain types have levels of at most $n-1$,
and that \emph{at least one} domain type has a level of exactly $n-1$.
This extension should still satisfy our desired properties.
Alternatively, if this structural change proves unviable,
multi-arity maps can be handled via an automatic front-end transformation to single-arity maps,
essentially treating the multi-arity syntax as syntactic sugar.


\subsection{No Map Assign Statement}
We currently omit support for Boogie's convenient map assignment syntax (e.g., \lstinline[columns=fixed]{M[x, y] := 1}), as discussed in Section~\ref{sec:semantics-boogie-maps}.
While this could be resolved via front-end desugaring into standard store operations,
extending the formal semantics to handle it natively would provide stronger correctness guarantees.

\subsection{Type Helper for Maps}
Certain components of the foundational Boogie semantics manage type inference and type safety.
These components must be explicitly extended to accommodate map cases.
Once this infrastructure is complete, the corresponding map cases can be implemented within the \isakeywordONE{ML} proof tactics that handle types.
This enhancement would allow a strictly larger subset of generated proofs to successfully verify.

\subsection{Current Automatic Proof Tactic Strength}
Currently, there is only rudimentary support for maps within the automated proof tactics.
We have included specific integer and boolean lemmas to verify common, simple use cases.
A deeper understanding of the existing \isakeywordONE{ML} tactics and extending them in a more general, exhaustive manner
is necessary to increase the verification success rate of generated proofs.

% \subsection{Non-trivial semantics}

% \subsection{Cleaning Up}
% % Map Semantics
% % Proof Gen Module
% The state of the map semantics file could be improved.
% Cleaner definition, and faster proofs.
% - simplify map semantics and proves -> preparation for extending more levels


\section{Alternative Approaches}

During the design phase, several alternative approaches were considered to address the complexities of modeling Boogie maps. We briefly outline these alternatives and the rationale behind our ultimate design choices.

\subsection{Association list and finite maps}

Our first idea was to model a map as an association list.
This is not a valid strategy, because Boogie maps are total functions and therefore infinite(see Section~\ref{sec:semantics-boogie-maps}).
Another idea, inspired by Section~2.1 of Burak Ekici et al.~\cite{ekici2017smtcoq},
is to use a finite mapping together with a default value for keys not contained in the mapping.
We argue that this is not appropriate for our use case,
because such a representation cannot capture all lambda expressions,
for example:

\begin{lstlisting}[numbers=none]
var m: [int]int;
m := (lambda x: int :: if i > 0 then 1 else -1);
\end{lstlisting}


\subsection{Semantics Based on Types Used in the Input Program}

One initial idea involved collecting all types used within the input Boogie program
and using this information to generate custom datatypes specific to the program's map types.
For example, if a Boogie program used an \lstinline[columns=fixed]{[int]int} map, we would model it directly as an \isa{int {\isasymRightarrow} int} function in Isabelle.
This stands in contrast to our current approach, which uniformly models all map levels using \isa{{\isacharprime}{\kern0pt}a\ val{\isadigit{0}\ {\isasymRightarrow}\ {\isacharprime}{\kern0pt}a\ val{\isadigit{0}}}} functions
(operating over all primitive values) and enforces type safety via a well-formedness invariant.

A significant hurdle with this approach was the representation of abstract types.
We later realized that Isabelle's \isakeywordONE{typedecl} command could potentially create types corresponding to these abstract types,
which could then be used to represent the domains and ranges of our maps.

The primary advantage of this approach is that it largely eliminates the need for complex well-formedness definitions.
The disadvantage, however, is the extensive overhead: one would need to statically extract all types, instantiate the map semantics for every observed type, and define polymorphic select and store operations for all possible type combinations. It remains an open question whether the complexity of generating these bespoke semantics outweighs the complexity of our unified, invariant-based approach.

\subsection{Context/Heap Model}

Another early concept was to maintain a separate global structure (a "heap") containing all maps,
accessed via a unique identifier (e.g., \isa{maps :: int {\isasymRightarrow} val {\isasymRightarrow} val}).
Under this model, a map value would merely store a reference, such as \isa{MapV int}.
This structure would have to be passed throughout the semantics and updated upon every map modification.

We discarded this approach because it introduces stateful complexity to the semantic model.
Reasoning about a heap structure introduces significant verification overhead.
Furthermore, determining the appropriate identifier format, and deciding if integers are suitable for quantifying over the infinite space of potential maps, posed a conceptual challenge.
Finally, establishing extensional equality between maps within a heap model is notably difficult.
While a context/heap model is excellent for implementing a runtime execution engine, it is poorly suited for a foundational logical model.

\subsection{Quotient Types}

A recurring issue was that our map datatypes were overly permissive, capable of representing values that have no valid Boogie counterpart
(e.g., non-type-safe maps or multiple representations of the same extensional Boogie map).
This motivated our rigorous well-formedness definition (Section~\ref{sec:wf-definition}).

Another potential avenue to enforce extensionality natively is the use of quotient types. \cite{huffman2013lifting}
Quotient types allow the definition of a new type by abstracting over multiple underlying representations using an equivalence relation.
Due to time constraints and the advanced Isabelle expertise required, we did not pursue this direction.
However, it represents a highly promising avenue for simplifying map semantics in the future.
For instance, one could model a map as an association list of updates with a fallback function,
and use quotienting to prove extensional equality between maps that were constructed differently but behave identically:

\begin{lstlisting}[caption={Two equivalent maps, built different}]
var a: [int]int;
var b: [int]int;
a := (lambda x:int :: x+1);
a[42] := 42;
b := (lambda x:int :: if x == 42 then 42 else x+1);

assert a == b;
\end{lstlisting}

\subsection{Deep Embedding via Lambda Calculus}

The most theoretically robust alternative is a deep embedding.
Our chosen solution is a shallow embedding, mapping Boogie maps directly to preexisting Isabelle functions.
Conversely, a deep embedding would model Boogie maps as expressions featuring explicit variable bindings.
Boogie maps closely mirror the lambda calculus:
Map selection (\lstinline[columns=fixed]{m[k]}) corresponds to function application (\isa{m k}),
map storage (\lstinline[columns=fixed]{m[k := v]}) corresponds to functional update (\isa{\isasymlambda x. if x = k then v else m x}),
and lambda expressions map directly to $\lambda$-abstractions.
With the inclusion of polymorphic map types (Section~\ref{sec:semantics-boogie-maps}), the type system closely resembles System F. % TODO ref
Therefore, the formal semantics of Boogie maps could be explicitly formalized as an implementation of System F.

This approach requires rigorously managing variables, environments, and De Bruijn indices or nominal bindings.
This significant conceptual overhead is why we initially pursued the simpler shallow embedding.
However, to formalize advanced features such as unbounded nesting and impredicative polymorphism, a deep embedding may ultimately be necessary.

A valid theoretical question is why a deep embedding avoids the paradoxes of Cantor's Theorem (Section~\ref{sec:cantor}).
The resolution lies in the fact that a deep embedding does not represent \emph{every} possible mathematical function in the HOL universe,
but only those computable functions expressible within the Boogie syntax.
For example, while the mathematical set of all integer functions (\lstinline[columns=fixed]{[int]bool}) is uncountable,
the set of functions constructible by a finite Boogie program is strictly countable.
\footnote{This assumes standard computable types. The inclusion of uncomputable reals could potentially complicate this, but within the bounds of a program, the expressible space remains constrained.}
By restricting the semantics to expressible functions, we bypass the cardinality paradoxes inherent in mapping full HOL function spaces to themselves.

\subsection{Alternative Logic Frameworks and Proof Assistants}

Given these challenges, one might ask whether Isabelle/HOL is the appropriate tool for defining map semantics.
Because the extensive prior work was built upon Isabelle/HOL, we elected to continue within this ecosystem.

However, if we wish to adjust the map semantics dynamically based on the nesting depth of the input program,
our datatype inherently depends on a numerical value.
This scenario strongly suggests the use of Dependent Types. %TODO ref
A dependently typed language (such as Coq, Lean, or Agda) could potentially define this hierarchy more elegantly.
Dependent types might also alleviate the need for our complex well-formedness predicates,
as constraints could be baked directly into the type signatures corresponding to the Boogie values.

% \subsection{Do we need one type for all values?}


\subsection{New Boogie Language: B3}

The Dafny development team is currently working on a new intermediate verification language: B3.\footnote{\url{https://github.com/dafny-lang/b3}}
If B3 ultimately supersedes the current Boogie IVL,
it raises the question of whether research efforts should pivot toward building foundational verification tools for this new ecosystem.
While B3 is still in its infancy, and the legacy Boogie IVL repository remains highly active,
it is an important development to monitor for future formalization projects.


\section{Future Work}

\subsection{Lambda Expressions}

One of the most compelling features of Boogie maps is the ability to assign lambda expressions to them (see Section~\ref{sec:semantics-boogie-maps}).
The current semantics have laid the groundwork, but substantial effort is still required to fully support this feature.
Naturally, one would first need to introduce a lambda expression datatype into the semantics.
The more complex task involves transforming this lambda expression into a function that fits the \isa{FunL} constructor,
and subsequently converting it into a valid, well-formed map value.
The primary challenge lies in formally proving that the resulting transformation satisfies all desired well-formedness properties,
while ensuring the intended functional behavior is accurately captured.
We have manually demonstrated the well-formedness of two specific lambda functions:

\begin{lstlisting}[numbers=none]
(lambda x : int :: x + 1) : [int]int
(lambda f : [int]int :: 
    (lambda x : int :: f[x] + 1) ) : [[int]int][int]int
\end{lstlisting}

Demonstrating the well-formedness of these examples highlights the complexity of dealing with stratification, sentinel values, and type safety.
Future implementations must devise a method to simplify this process and automate both the generation and the corresponding proofs.
This is a highly promising avenue for future work.

\subsection{Polymorphic Map Types}

Arguably, the most complex feature of Boogie maps is their support for polymorphic types (see Section~\ref{sec:semantics-boogie-maps}),
and the resulting impredicative polymorphism and recursive types.
It even permits the assignment of polymorphic lambdas.

This functionality completely breaks our stratification approach (Section~\ref{sec:stratification-solution}),
as a map of level $n$ fundamentally cannot possess a domain of level $n$.
It remains unclear under which circumstances such highly cyclic maps remain logically consistent.
Supporting such maps would require an entirely different semantic approach.
However, as noted in prior work, such pathological use cases rarely appear in practice,
suggesting that one might safely restrict the generality of map semantics to accommodate standard, practical verification tasks.

\subsection{Monomorphic Type Encoding}

The entire proof generation module currently relies on Boogie's predicate type encoding.
At the time of the prior work, this was the default encoding;
however, the default has since switched to monomorphic type encoding.
Adapting the proof generation to support monomorphic type encoding could significantly benefit map semantics.
The monomorphic encoding uses more restrictive SMT Array Theory axioms in the verification conditions (Section~\ref{sec:type-encoding}).
Leveraging these restrictions could facilitate a simpler formulation of the map semantics, potentially enabling more general map behaviors.

\subsection{Addressing Current Limitations}

A natural direction for future work is addressing the limitations outlined in Section~\ref{sec:limitations}.
This includes natively supporting multi-arity maps, implementing the convenient map assignment syntax,
formally verifying map-related type safety lemmas,
and strengthening the \isakeywordONE{ML} proof tactics to enable a higher rate of automated proof verification.

\subsection{Extending Map Nesting Levels}

Our solution currently supports maps with up to three levels of nesting in their domain types.
There is no theoretical limit for nesting range types, as recursion on the range is logically safe (see Section~\ref{sec:cantor}).
Because deeply nested domains are rare in practical benchmarks, a three-level limit serves as a highly practical bound.
Nevertheless, a fully generalized semantics is desirable.
Extending our approach to support an arbitrarily large, bounded (or even unbounded) number of nesting levels is a clear next step.
One approach is to programmatically determine the required nesting level from the input program
and dynamically generate the corresponding Isabelle semantics, definitions, and proofs.
An even more robust solution would involve defining a single semantic framework capable of handling arbitrary depths without relying on program-dependent generation.

\subsection{Supporting Additional Boogie Features}

A primary motivation for this work (Section~\ref{sec:motivation}) is to enable a foundational verification pipeline capable of processing Dafny programs.
While we resolved a significant hurdle by formalizing map semantics, other unsupported Boogie features currently prevent the verification of even simple Dafny examples.
For instance, the proof generation module currently fails when encountering procedures marked with \lstinline[columns=fixed]{free} postconditions.
Furthermore, several value conversion functions (e.g., converting a \lstinline[columns=fixed]{real} to an \lstinline[columns=fixed]{int}) remain unimplemented.
Addressing these unsupported features is the necessary next step toward achieving fully foundational verification for Dafny programs.

% \subsection{Improve Proofgen code}
% - new CLI interface
% - reduce intrusion
